\chapter{استقلال خطی}
\label{ch:independence}

وقتی زیرفضایی را با چند بردار تولید می‌کنیم، ممکن است بعضی بردارها \emph{اضافه} باشند؛ یعنی خودشان ترکیبی از بقیه باشند. \textbf{استقلال خطی} همین «اضافه نبودن» را دقیق می‌کند و راه را برای \emph{پایه} و \emph{بُعد} در فصل بعد باز می‌کند.

\section{انگیزه: قیمت خانه}

می‌خواهیم قیمت خانه ($y$) را از متراژ ($f$) با مدل $y=af+c$ پیش‌بینی کنیم. هر داده‌ی $(f^{(i)},y^{(i)})$ یک معادله‌ی $af^{(i)}+c=y^{(i)}$ می‌دهد و همه با هم دستگاه
\[
\begin{bmatrix}f^{(1)}&1\\ \vdots&\vdots\\ f^{(N)}&1\end{bmatrix}\begin{bmatrix}a\\c\end{bmatrix}=\begin{bmatrix}y^{(1)}\\ \vdots\\ y^{(N)}\end{bmatrix}
\]
را می‌سازند. با ویژگی‌های بیشتر (تعداد اتاق، سن بنا، متراژ بخش‌ها و …) هر ویژگی یک \textbf{ستون} $A$ می‌شود. اگر «متراژ کل = مجموع متراژ بخش‌ها» باشد، آن ستون ترکیب خطی بقیه است و هیچ اطلاع تازه‌ای نمی‌دهد. زبان دقیق این بحث، استقلال خطی است.

\section{تعریف}

\begin{definition}[وابستگی و استقلال خطی]
بردارهای $\vv_1,\dots,\vv_n$ \term{مستقل خطی}{Linearly Independent} هستند اگر
\[
\lambda_1\vv_1+\dots+\lambda_n\vv_n=\zero\ \Longrightarrow\ \lambda_1=\dots=\lambda_n=0 .
\]
در غیر این صورت \term{وابسته‌ی خطی}{Linearly Dependent} هستند؛ یعنی رابطه‌ای با دست‌کم یک $\lambda_i\neq0$ وجود دارد (\textbf{رابطه‌ی وابستگی}).
\end{definition}

به زبان ساده (قضیه‌ی \ref{thm:dep-char}): مجموعه وابسته است اگر دست‌کم یکی از بردارهایش ترکیب خطی بقیه باشد.

\begin{keyformula}[top=2mm,bottom=2mm]
\centering
ستون‌های $A$ مستقل‌اند
\ $\iff$\
$A\vx=\zero$ فقط جواب بدیهی دارد
\ $\iff$\
همه‌ی ستون‌ها محوری‌اند.
\end{keyformula}
(تساوی دوم از قضیه‌ی \ref{thm:homog-free} می‌آید.)

\begin{intuition}
دو بردار وابسته‌اند اگر و تنها اگر \textbf{هم‌راستا} باشند؛ سه بردار در $\R^3$ وابسته‌اند اگر و تنها اگر \textbf{هم‌صفحه} باشند.
\begin{center}
\begin{tikzpicture}[scale=0.8,>=Latex]
  \draw[->, very thick] (0,0)--(-0.9,2.2);
  \draw[->, very thick, gray] (0,0)--(-0.45,1.1);
  \node at (-0.2,-0.4) {\small \fa{وابسته}};
  \begin{scope}[xshift=3cm]
  \draw[->, very thick] (0,0)--(-0.9,2.2);
  \draw[->, very thick, gray] (0,0)--(0.7,1.6);
  \node at (-0.1,-0.4) {\small \fa{مستقل}};
  \end{scope}
  \begin{scope}[xshift=6cm]
  \draw[->, very thick] (0,0)--(-0.9,2.2);
  \draw[->, very thick] (0,0)--(0.1,1.9);
  \draw[->, very thick, gray] (0,0)--(0.7,1.6);
  \node at (-0.1,-0.4) {\small \fa{وابسته (سه بردار در $\R^2$)}};
  \end{scope}
\end{tikzpicture}
\end{center}
\end{intuition}

\begin{example}
آیا $\vv_1=(1,2,3)$، $\vv_2=(4,5,6)$ و $\vv_3=(2,1,0)$ مستقل‌اند؟
\end{example}
\begin{solution}
\[
\begin{bmatrix}1&4&2\\2&5&1\\3&6&0\end{bmatrix}\sim\begin{bmatrix}1&0&-2\\0&1&1\\0&0&0\end{bmatrix}.
\]
ستون سوم محوری نیست؛ پس وابسته‌اند. با $x_3=1$: $x_1=2,\ x_2=-1$، یعنی $2\vv_1-\vv_2+\vv_3=\zero$.
\end{solution}

\section{ویژگی‌ها}

\begin{theorem}\label{thm:zero-dep}
هر مجموعه‌ی شامل $\zero$ وابسته است.
\end{theorem}
\begin{proof}
در $0\vv_1+\dots+0\vv_k+1\cdot\zero=\zero$ ضریب $\zero$ ناصفر است.
\end{proof}

\begin{theorem}[مشخصه‌ی مجموعه‌های وابسته]\label{thm:dep-char}
مجموعه‌ی $\{\vv_1,\dots,\vv_n\}$ ($n\ge2$) وابسته است اگر و تنها اگر یکی از بردارها ترکیب خطی بقیه باشد. در واقع اگر وابسته باشد و $\vv_1\neq\zero$، یک $\vv_j$ ($j>1$) ترکیب خطی \textbf{بردارهای قبل از خودش} است.
\end{theorem}
\begin{proof}
\textbf{(برگشت)} اگر $\vv_j=\sum_{i\neq j}c_i\vv_i$، آن‌گاه $\sum_{i\neq j}c_i\vv_i-\vv_j=\zero$ رابطه‌ای با ضریب ناصفر $-1$ است.

\textbf{(رفت)} اگر $\vv_1=\zero$، $\vv_1=0\vv_2+\dots+0\vv_n$. وگرنه رابطه‌ی $c_1\vv_1+\dots+c_n\vv_n=\zero$ با ضرایب نه‌همه‌صفر داریم. فرض کنید $j$ بزرگ‌ترین اندیس با $c_j\neq0$ باشد. اگر $j=1$ بود، $c_1\vv_1=\zero$ و طبق قضیه‌ی \ref{thm:vs-basic} $\vv_1=\zero$؛ ناممکن. پس $j>1$ و
\[
\vv_j=-\frac{c_1}{c_j}\vv_1-\dots-\frac{c_{j-1}}{c_j}\vv_{j-1}.
\]
\end{proof}

\begin{warning}
این قضیه نمی‌گوید \emph{هر} بردار ترکیب خطی بقیه است: در $\{(1,0),(2,0),(0,1)\}$ بردار $(0,1)$ ترکیب خطی بقیه نیست، ولی مجموعه وابسته است.
\end{warning}

\begin{theorem}\label{thm:more-than-n}
هر $p$ بردار در $\R^n$ با $p>n$ وابسته‌اند.
\end{theorem}
\begin{proof}
$A\vx=\zero$ با $A=\begin{bmatrix}\vv_1&\cdots&\vv_p\end{bmatrix}$ دستگاهی همگن با $n$ معادله و $p>n$ مجهول است و طبق نتیجه‌ی \ref{cor:more-unknowns} جواب نابدیهی دارد.
\end{proof}

\begin{example}
وابسته یا مستقل؟ (بدون محاسبه)
(الف) چهار بردار در $\R^3$؛
(ب) $(2,3,5)$، $(0,0,0)$، $(1,1,8)$؛
(ج) $(-2,4,6,10)$ و $(3,-6,-9,15)$.
\end{example}
\begin{solution}
(الف) وابسته (قضیه‌ی \ref{thm:more-than-n}). (ب) وابسته (قضیه‌ی \ref{thm:zero-dep}). (ج) مستقل: از مؤلفه‌ی اول نسبت $-\frac32$ و از مؤلفه‌ی آخر $\frac32$ به دست می‌آید؛ پس مضرب هم نیستند.
\end{solution}

\begin{theorem}[حذف بردار اضافی]\label{thm:remove-dependent}
اگر $c_1\vv_1+\dots+c_n\vv_n=\zero$ و $c_1\neq0$، آن‌گاه
$\Span\{\vv_1,\dots,\vv_n\}=\Span\{\vv_2,\dots,\vv_n\}$.
\end{theorem}
\begin{proof}
$\vv_1=-\frac{c_2}{c_1}\vv_2-\dots-\frac{c_n}{c_1}\vv_n$؛ پس در هر ترکیب خطی می‌شود $\vv_1$ را جایگذاری کرد و ترکیبی از $\vv_2,\dots,\vv_n$ به دست آورد.
\end{proof}

با تکرار این حذف به فهرستی مستقل می‌رسیم که همان فضا را تولید می‌کند: یک \emph{پایه} (فصل ۶).

\begin{theorem}[یکتایی نمایش]\label{thm:unique-coeffs}
اگر $\vv_1,\dots,\vv_n$ مستقل باشند، ضرایب نمایش $\vx=\beta_1\vv_1+\dots+\beta_n\vv_n$ یکتا هستند.
\end{theorem}
\begin{proof}
اگر $\vx=\sum\beta_i\vv_i=\sum\gamma_i\vv_i$، آن‌گاه $\sum(\beta_i-\gamma_i)\vv_i=\zero$ و از استقلال $\beta_i=\gamma_i$.
\end{proof}

\begin{theorem}\label{thm:subset-superset}
ابرمجموعه‌ی هر مجموعه‌ی وابسته، وابسته است؛ و زیرمجموعه‌ی ناتهی هر مجموعه‌ی مستقل، مستقل است.
\end{theorem}
\begin{proof}
به رابطه‌ی وابستگی، بردارهای جدید را با ضریب صفر اضافه کنید. بخش دوم عکس نقیض بخش اول است.
\end{proof}

\begin{theorem}\label{thm:null-vectors-indep}
بردارهای شکل برداری جواب عمومی $A\vx=\zero$ (یکی برای هر متغیر آزاد) مستقل‌اند.
\end{theorem}
\begin{proof}
اگر $\vx=x_{f_1}\vu_1+\dots+x_{f_k}\vu_k$، مؤلفه‌ی $f_i$ام $\vu_i$ برابر ۱ و همین مؤلفه در بقیه‌ی $\vu_j$ها صفر است. پس در $\sum c_j\vu_j=\zero$، مؤلفه‌ی $f_i$ام می‌دهد $c_i=0$.
\end{proof}

\begin{note}[روش بررسی $p$ بردار در $\R^n$]
(۱) اگر $p>n$ وابسته‌اند. (۲) اگر $\zero$ در میان است وابسته‌اند. (۳) دو بردار وابسته‌اند اگر مضرب هم باشند. (۴) در حالت کلی، کاهش سطری: مستقل‌اند اگر همه‌ی ستون‌ها محوری باشند.
\end{note}

\section{استقلال توابع و چندجمله‌ای‌ها}

در فضای توابع، «$=0$» یعنی برای \textbf{همه‌ی} $t$ها. توابع $f,g$ وابسته‌اند اگر $c_1,c_2$ نه‌هر‌دو‌صفر باشند با $c_1f(t)+c_2g(t)=0$ برای همه‌ی $t$.

\begin{example}
(الف) $2\sin^2t$ و $1-\cos^2t$؛ (ب) $\sin^2x$، $\cos^2x$ و $\cos2x$؛ (ج) $\sin t$ و $\cos t$.
\end{example}
\begin{solution}
(الف) \textbf{وابسته:} $2\sin^2t-2(1-\cos^2t)=0$. \
(ب) \textbf{وابسته:} $\sin^2x-\cos^2x+\cos2x=0$. \
(ج) \textbf{مستقل:} اگر $c_1\sin t+c_2\cos t=0$ برای همه‌ی $t$، با $t=0$ داریم $c_2=0$ و با $t=\frac\pi2$ داریم $c_1=0$.
\end{solution}

\begin{warning}
برای \emph{وابستگی} باید رابطه را برای همه‌ی $t$ها نشان دهید (با یک اتحاد)؛ برای \emph{استقلال} کافی است چند مقدار $t$ ضرایب را صفر کنند.
\end{warning}

\begin{example}
آیا $1-x$، $1+x$ و $x^2$ مستقل‌اند؟
\end{example}
\begin{solution}
$c_1(1-x)+c_2(1+x)+c_3x^2=(c_1+c_2)+(c_2-c_1)x+c_3x^2=0$ یعنی همه‌ی ضرایب صفرند: $c_1+c_2=0$، $c_2-c_1=0$، $c_3=0$. پس $c_1=c_2=c_3=0$ و مستقل‌اند.
\end{solution}

\begin{theorem}\label{thm:monomials-indep}
$1,x,\dots,x^n$ مستقل خطی‌اند.
\end{theorem}
\begin{proof}
اگر $c_0+c_1x+\dots+c_nx^n=0$ برای همه‌ی $x$ و ضرایب همه صفر نباشند، یک چندجمله‌ای ناصفر با درجه‌ی حداکثر $n$ بی‌شمار ریشه دارد؛ تناقض.
\end{proof}

\section{جمع‌بندی و تمرین}

\begin{summary}
\begin{itemize}
  \item مستقل: $\sum\lambda_i\vv_i=\zero$ فقط جواب بدیهی دارد؛ یعنی هیچ برداری ترکیب خطی بقیه نیست.
  \item ستون‌های $A$ مستقل‌اند $\iff$ همه‌ی ستون‌ها محوری‌اند.
  \item شامل $\zero$ یا بیش از $n$ بردار در $\R^n$ $\Rightarrow$ وابسته.
  \item بردار وابسته را می‌شود بدون تغییر پوشش حذف کرد؛ نمایش بر حسب بردارهای مستقل یکتاست.
\end{itemize}
\end{summary}

\subsection*{تمرین‌ها}
\begin{enumerate}
  \item به ازای چه $h$ بردارهای $(1,-1,4)$، $(3,-5,7)$ و $(-1,5,h)$ وابسته‌اند؟
  \item ثابت کنید اگر $\{\vu,\vv\}$ مستقل باشد، $\{\vu+\vv,\vu-\vv\}$ هم مستقل است.
  \item درست یا نادرست: «اگر $\vv_3$ ترکیب خطی $\vv_1,\vv_2$ نباشد، $\{\vv_1,\vv_2,\vv_3\}$ مستقل است.»
\end{enumerate}

\subsection*{پاسخ‌های کوتاه}
\begin{enumerate}
  \item درایه‌ی محوری سوم $h-6$ می‌شود؛ پس فقط برای $h=6$.
  \item $a(\vu+\vv)+b(\vu-\vv)=(a+b)\vu+(a-b)\vv=\zero$؛ پس $a+b=a-b=0$ و $a=b=0$.
  \item نادرست: $(1,0)$، $(2,0)$، $(0,1)$.
\end{enumerate}
